\documentclass[a4paper,12pt]{article}
\usepackage[ngerman]{babel}
\usepackage[utf8]{inputenc}
\usepackage[T1]{fontenc}
\usepackage{amsmath,amssymb}
\usepackage{tikz}
\usetikzlibrary{positioning,decorations.pathreplacing}

\newcommand{\methodbox}[1]{%
\medskip
\noindent\fbox{%
    \begin{minipage}{0.94\linewidth}
    \textbf{Methodik / Definition.}\par\smallskip
    \small
    #1
    \end{minipage}%
}
\medskip\par
}

\newcommand{\block}[4]{%
    \draw[fill=#4!25] (#1,0) rectangle ++(#2,0.65);
    \node at ({#1 + #2/2},0.325) {#3};
}

\begin{document}
\raggedright
\setlength{\parindent}{0pt}

{\LARGE\bfseries Pumping-Lemma: einfaches Beispiel\par}
\vspace{1.5em}

\methodbox{
Pumping-Lemma zum Nachweis von Nicht-Regularität:
\begin{enumerate}
    \item Annahme: Die Sprache ist regulär. Sei $p$ die Pumping-Länge.
    \item Wähle ein Wort $s$ aus der Sprache mit $|s|\ge p$.
    \item Betrachte eine beliebige Zerlegung $s=xyz$ mit $|y|>0$ und $|xy|\le p$.
    \item Zeige für ein geeignetes $i$ meistens $i=0$ oder $i=2$, dass $xy^iz$ nicht mehr in der Sprache liegt.
\end{enumerate}
Das ist ein Widerspruch. Also ist die Sprache nicht regulär.
}

Wir betrachten die Sprache
\[
L=\{w\in\{a,b\}^*\mid |w|_a>|w|_b\}.
\]
Dabei bezeichnet $|w|_a$ die Anzahl der $a$'s in $w$ und $|w|_b$ die Anzahl der $b$'s in $w$.

\medskip
\textbf{Beweis.}
\begin{enumerate}
    \item Annahme: $L$ ist regulär. Sei $p$ die Pumping-Länge.

    \item Wähle
    \[
        s=a^{p+1}b^p.
    \]
    Dann gilt
    \[
        |s|_a=p+1>p=|s|_b,
    \]
    also ist $s\in L$. Außerdem gilt $|s|=2p+1\ge p$.

    \begin{center}
    \begin{tikzpicture}[x=0.9cm,y=0.9cm]
        \block{0}{5}{$a^{p+1}$}{blue}
        \block{5}{4}{$b^p$}{red}
        \draw[decorate,decoration={brace,mirror,amplitude=5pt}] (0,-0.1) -- (5,-0.1)
            node[midway,below=7pt] {$p+1$ viele $a$'s};
        \draw[decorate,decoration={brace,mirror,amplitude=5pt}] (5,-0.1) -- (9,-0.1)
            node[midway,below=7pt] {$p$ viele $b$'s};
    \end{tikzpicture}
    \end{center}

    \item Sei $s=xyz$ eine Zerlegung gemäß Pumping-Lemma, also
    \[
        |y|>0
        \qquad\text{und}\qquad
        |xy|\le p.
    \]
    Weil die ersten $p$ Zeichen von $s$ alle $a$ sind, liegt $y$ vollständig im ersten $a$-Block.
    Also gilt
    \[
        y=a^k
        \qquad\text{für ein } k\ge 1.
    \]

    \begin{center}
    \begin{tikzpicture}[x=0.9cm,y=0.9cm]
        \block{0}{1.6}{$x$}{gray}
        \block{1.6}{1.4}{$y=a^k$}{green}
        \block{3}{2}{$a^{p+1-k-|x|}$}{blue}
        \block{5}{4}{$b^p$}{red}
        \draw[decorate,decoration={brace,mirror,amplitude=5pt}] (0,-0.1) -- (5,-0.1)
            node[midway,below=7pt] {nur $a$'s};
    \end{tikzpicture}
    \end{center}

    \item Wir pumpen mit $i=0$, also entfernen wir $y$:
    \[
        xy^0z=xz=a^{p+1-k}b^p.
    \]

    \item Jetzt gilt
    \[
        |xz|_a=p+1-k
        \qquad\text{und}\qquad
        |xz|_b=p.
    \]
    Wegen $k\ge 1$ folgt
    \[
        p+1-k\le p.
    \]
    Also hat $xz$ nicht mehr mehr $a$'s als $b$'s:
    \[
        xz\notin L.
    \]

    \begin{center}
    \begin{tikzpicture}[x=0.9cm,y=0.9cm]
        \block{0}{4}{$a^{p+1-k}$}{blue}
        \block{4}{4}{$b^p$}{red}
        \node[below=0.65cm] at (2,0) {$p+1-k\le p$ viele $a$'s};
        \node[below=0.65cm] at (6,0) {$p$ viele $b$'s};
    \end{tikzpicture}
    \end{center}

    Das ist ein Widerspruch zum Pumping-Lemma.
\end{enumerate}

Folglich ist $L$ nicht regulär.

\end{document}
